\clearpage
\section*{ПРИЛОЖЕНИЕ А}
\addcontentsline{toc}{section}{ПРИЛОЖЕНИЕ А}

\begin{flushright}
\textbf{Приложение А}
\end{flushright}

\begin{center}
\textbf{СТРУКТУРА ПРОГРАММНОГО ПРОЕКТА И ПРИМЕРЫ ЗАПУСКА}
\end{center}

В приложении приведена укрупненная структура программного проекта, разработанного в рамках настоящей выпускной квалификационной работы.

\begin{verbatim}
diploma_rl_trading/
├── data/
├── notebooks/
├── src/
│   ├── data/
│   ├── strategies/
│   ├── models/
│   ├── rl/
│   ├── metrics/
│   └── visualization/
├── reports/
├── requirements.txt
└── main.py
\end{verbatim}

Назначение ключевых элементов структуры:
\begin{enumerate}[leftmargin=1.25cm,label=\arabic*)]
\item каталог \texttt{data/} используется для хранения исходных и обработанных котировок;
\item каталог \texttt{notebooks/} предназначен для исследовательских записных книжек и промежуточного анализа;
\item каталог \texttt{src/} содержит основную логику системы, разбитую по функциональным модулям;
\item каталог \texttt{reports/} служит для сохранения метрик, таблиц, графиков и прочих результатов экспериментов;
\item файл \texttt{requirements.txt} фиксирует зависимости проекта;
\item файл \texttt{main.py} запускает полный сравнительный pipeline.
\end{enumerate}

Ниже приведены примеры команд, использованных в рамках экспериментального моделирования:

\begin{verbatim}
python3 main.py --ticker AAPL --start_date 2018-01-01 --end_date 2024-01-01 \
  --initial_cash 10000 --transaction_cost 0.001 --random_state 42

python3 src/rl/train_ppo.py --data_path data/processed/aapl_features.csv \
  --ticker AAPL --timesteps 20000 --initial_cash 10000 --transaction_cost 0.001

python3 src/rl/train_dqn.py --data_path data/processed/aapl_features.csv \
  --ticker AAPL --timesteps 15000 --initial_cash 10000 --transaction_cost 0.001
\end{verbatim}

Представленная структура проекта обеспечивает воспроизводимость вычислительного эксперимента и позволяет последовательно расширять стенд за счет новых признаков, алгоритмов и сценариев оценки.
